\documentclass[a4paper,12pt]{article}
\usepackage[utf8]{inputenc}
\usepackage[ngerman]{babel}
\usepackage{amsmath,amssymb}
\usepackage{geometry}\geometry{margin=2.5cm}
\usepackage{enumitem}
\usepackage{booktabs}
\usepackage{tcolorbox}
\tcbuselibrary{breakable}
\usepackage{xcolor}

\newtcolorbox{begriffe}[1][]{colback=yellow!10, colframe=yellow!60!black, fonttitle=\bfseries, title={Was bedeuten die Begriffe?}, breakable, #1}
\newtcolorbox{wastun}[1][]{colback=orange!5, colframe=orange!60!black, fonttitle=\bfseries, title={Was ist zu tun?}, breakable, #1}
\newcommand{\piv}[1]{{\color{red}\boxed{#1}}}

\title{\"Ubungsblatt 6 -- L\"osungen \\ \large Linear Algebra}
\author{}
\date{}

\begin{document}
\maketitle

\begin{center}
\fbox{\parbox{0.9\textwidth}{
Die Abschnitte mit $\square$ sind genau die einzelnen Beispiele zum Kreuzen:
22(a), 22(b), 22(c), 23(a), 23(b) und 24.
}}
\end{center}

%% ============================================================
\subsection*{$\square$ Beispiel 22(a)}
%% ============================================================

\textbf{Angabe.}
\[
W = \{(a,b,c,d,e)\in \mathbb{R}^5 : a+b+c = a+2b-3d = a+3b-5e\}.
\]
Gesucht sind eine Basis von $W$ und die Dimension von $W$.

\begin{begriffe}
\begin{itemize}[leftmargin=*, nosep]
    \item $W$ ist eine Menge von Vektoren in $\mathbb{R}^5$.
    \item Eine \textbf{Basis} ist eine Liste von Vektoren, mit denen man alle Vektoren in $W$ bauen kann, ohne dass einer davon \"uberfl\"ussig ist.
    \item Die \textbf{Dimension} ist die Anzahl der Vektoren in einer Basis.
    \item Die Schreibweise $T_1=T_2=T_3$ bedeutet: alle drei Ausdr\"ucke sind gleich. Daf\"ur reicht es, $T_1=T_2$ und $T_2=T_3$ zu fordern.
\end{itemize}
\end{begriffe}

\begin{wastun}
\begin{enumerate}[leftmargin=*, nosep]
    \item Aus der Kettengleichung zwei normale Gleichungen machen.
    \item Diese Gleichungen nach einigen Variablen aufl\"osen.
    \item Die freien Variablen als Parameter schreiben.
    \item Aus der Parameterdarstellung die Basis ablesen.
\end{enumerate}
\end{wastun}

\textbf{Schritt 1: Die Kettengleichung in zwei Gleichungen zerlegen.}

Wir setzen
\[
T_1=a+b+c,\qquad T_2=a+2b-3d,\qquad T_3=a+3b-5e.
\]
Die Bedingung $T_1=T_2=T_3$ ist gleichbedeutend mit
\[
T_1=T_2
\qquad\text{und}\qquad
T_2=T_3.
\]

Erste Gleichung:
\begin{align*}
a+b+c &= a+2b-3d \\
b+c &= 2b-3d && \text{auf beiden Seiten $a$ abziehen} \\
c-b+3d &= 0 && \text{alles auf eine Seite bringen} \\
c &= b-3d.
\end{align*}

Zweite Gleichung:
\begin{align*}
a+2b-3d &= a+3b-5e \\
2b-3d &= 3b-5e && \text{auf beiden Seiten $a$ abziehen} \\
-b-3d+5e &= 0 && \text{alles auf eine Seite bringen} \\
5e &= b+3d \\
e &= \frac{b+3d}{5}.
\end{align*}

\textbf{Schritt 2: Freie Variablen erkennen.}

Die Variable $a$ kommt in den Endbedingungen gar nicht mehr vor. Also ist $a$ frei.
Au\ss erdem d\"urfen wir $b$ und $d$ frei w\"ahlen. Danach sind $c$ und $e$ festgelegt:
\[
c=b-3d,
\qquad
e=\frac{b+3d}{5}.
\]

Damit sieht jeder Vektor aus $W$ so aus:
\[
(a,b,c,d,e)
=
\left(a,\; b,\; b-3d,\; d,\; \frac{b+3d}{5}\right).
\]

\textbf{Schritt 3: Parameter so w\"ahlen, dass keine Br\"uche in der Basis stehen.}

Wir setzen
\[
a=r,\qquad b=5s,\qquad d=5t
\quad (r,s,t\in\mathbb{R}).
\]
Das ist keine Einschr\"ankung: Weil $s$ und $t$ beliebige reelle Zahlen sind, k\"onnen $5s$ und $5t$ immer noch jeden reellen Wert annehmen.

Dann ist
\[
c = 5s-15t,
\qquad
e = \frac{5s+15t}{5}=s+3t.
\]
Also
\begin{align*}
(a,b,c,d,e)
&= (r,\;5s,\;5s-15t,\;5t,\;s+3t) \\
&= r(1,0,0,0,0)
 + s(0,5,5,0,1)
 + t(0,0,-15,5,3).
\end{align*}

\textbf{Schritt 4: Basis und Dimension ablesen.}

Die drei Parameter $r,s,t$ sind frei. Zu jedem freien Parameter geh\"ort ein Basisvektor.

\[
\fbox{\parbox{0.93\textwidth}{\centering
Eine Basis von $W$ ist
\[
\{(1,0,0,0,0),\; (0,5,5,0,1),\; (0,0,-15,5,3)\}.
\]
Also ist $\dim(W)=3$.
}}
\]

\textbf{Kurze Kontrolle.}
Wenn man die drei Basisvektoren in die urspr\"ungliche Bedingung einsetzt, erf\"ullen sie jeweils
$a+b+c=a+2b-3d=a+3b-5e$. Au\ss erdem kann jeder Vektor in $W$ mit den drei Parametern $r,s,t$ gebaut werden. Daher ist das wirklich eine Basis.

\bigskip
\hrule
\bigskip

%% ============================================================
\subsection*{$\square$ Beispiel 22(b)}
%% ============================================================

\textbf{Angabe.}
$W$ ist der Raum aller reellen Matrizen der Form
\[
A=
\begin{pmatrix}
a & b & c \\
b & c & d \\
c & d & e
\end{pmatrix}
\]
mit
\[
A
\begin{pmatrix}
1\\-2\\3
\end{pmatrix}
=
\begin{pmatrix}
0\\0\\0
\end{pmatrix}.
\]
Gesucht sind eine Basis von $W$ und die Dimension von $W$.

\begin{begriffe}
\begin{itemize}[leftmargin=*, nosep]
    \item Die Matrix ist nicht beliebig: Sie ist durch die f\"unf Zahlen $a,b,c,d,e$ bestimmt.
    \item Die Bedingung $A(1,-2,3)^T=0$ liefert drei Gleichungen.
    \item Wieder suchen wir alle m\"oglichen Matrizen, die diese Gleichungen erf\"ullen.
\end{itemize}
\end{begriffe}

\textbf{Schritt 1: Matrix-Vektor-Produkt ausrechnen.}

\[
\begin{pmatrix}
a & b & c \\
b & c & d \\
c & d & e
\end{pmatrix}
\begin{pmatrix}
1\\-2\\3
\end{pmatrix}
=
\begin{pmatrix}
a-2b+3c \\
b-2c+3d \\
c-2d+3e
\end{pmatrix}.
\]

Dieser Vektor soll der Nullvektor sein. Also m\"ussen alle drei Eintr\"age gleich $0$ sein:
\[
\begin{cases}
a-2b+3c=0,\\
b-2c+3d=0,\\
c-2d+3e=0.
\end{cases}
\]

\textbf{Schritt 2: Von unten nach oben aufl\"osen.}

Die letzte Gleichung enth\"alt $c,d,e$:
\begin{align*}
c-2d+3e &= 0 \\
c &= 2d-3e.
\end{align*}

Jetzt setzen wir das in die zweite Gleichung ein:
\begin{align*}
b-2c+3d &= 0 \\
b &= 2c-3d \\
&= 2(2d-3e)-3d \\
&= 4d-6e-3d \\
&= d-6e.
\end{align*}

Jetzt setzen wir $b=d-6e$ und $c=2d-3e$ in die erste Gleichung ein:
\begin{align*}
a-2b+3c &= 0 \\
a &= 2b-3c \\
&= 2(d-6e)-3(2d-3e) \\
&= 2d-12e-6d+9e \\
&= -4d-3e.
\end{align*}

\textbf{Fertige Gleichungen f\"ur alle Buchstaben.}

Nach Schritt 2 wissen wir also:
\[
\begin{array}{rcl}
a &=& -4d-3e,\\
b &=& d-6e,\\
c &=& 2d-3e,\\
d &=& d,\\
e &=& e.
\end{array}
\]

Das bedeutet: Sobald wir $d$ und $e$ frei w\"ahlen, sind $a$, $b$ und $c$ automatisch festgelegt.
Es bleiben also genau zwei freie Zahlen \"ubrig; daraus werden gleich die zwei Parameter f\"ur die Basis.

\textbf{Schritt 3: Freie Variablen setzen.}

Die Variablen $d$ und $e$ sind frei. Setze
\[
d=s,\qquad e=t
\quad (s,t\in\mathbb{R}).
\]
Dann gilt
\[
c=2s-3t,\qquad b=s-6t,\qquad a=-4s-3t.
\]

Damit ist
\[
A=
\begin{pmatrix}
-4s-3t & s-6t & 2s-3t \\
s-6t & 2s-3t & s \\
2s-3t & s & t
\end{pmatrix}.
\]

Jetzt trennen wir die $s$-Anteile und die $t$-Anteile:
\[
A
=
s
\begin{pmatrix}
-4 & 1 & 2 \\
1 & 2 & 1 \\
2 & 1 & 0
\end{pmatrix}
+
t
\begin{pmatrix}
-3 & -6 & -3 \\
-6 & -3 & 0 \\
-3 & 0 & 1
\end{pmatrix}.
\]

\textbf{Schritt 4: Basis und Dimension ablesen.}

\[
\fbox{\parbox{0.93\textwidth}{\centering
Eine Basis von $W$ ist
\[
\left\{
\begin{pmatrix}
-4 & 1 & 2 \\
1 & 2 & 1 \\
2 & 1 & 0
\end{pmatrix},
\begin{pmatrix}
-3 & -6 & -3 \\
-6 & -3 & 0 \\
-3 & 0 & 1
\end{pmatrix}
\right\}.
\]
Also ist $\dim(W)=2$.
}}
\]

\textbf{Warum sind es genau zwei Basisvektoren?}
Es gibt genau zwei freie Parameter, n\"amlich $s$ und $t$.
Jede erlaubte Matrix entsteht aus diesen zwei Parametern.
Die beiden Basisvektoren sind auch nicht Vielfache voneinander, zum Beispiel hat der erste unten rechts den Eintrag $0$, der zweite unten rechts den Eintrag $1$.

\bigskip
\hrule
\bigskip

%% ============================================================
\subsection*{$\square$ Beispiel 22(c)}
%% ============================================================

\textbf{Angabe.}
\[
W=\{p(x)\in \mathbb{R}_4[x] : p(0)=p(1)=p(2)\}.
\]
Gesucht sind eine Basis von $W$ und die Dimension von $W$.

\begin{begriffe}
\begin{itemize}[leftmargin=*, nosep]
    \item $\mathbb{R}_4[x]$ bedeutet: alle Polynome mit Grad h\"ochstens $4$.
    \item Die Bedingung $p(0)=p(1)=p(2)$ sagt: Bei $x=0$, $x=1$ und $x=2$ hat das Polynom denselben Wert.
    \item Wenn ein Polynom $q$ bei $x=r$ den Wert $0$ hat, dann ist $(x-r)$ ein Faktor von $q$.
\end{itemize}
\end{begriffe}

\textbf{Schritt 1: Den gemeinsamen Wert abziehen.}

Setze
\[
\lambda = p(0).
\]
Weil $p(0)=p(1)=p(2)$ gilt, ist auch
\[
p(0)=\lambda,\qquad p(1)=\lambda,\qquad p(2)=\lambda.
\]

Jetzt betrachten wir
\[
q(x)=p(x)-\lambda.
\]
Dann gilt:
\[
q(0)=p(0)-\lambda=0,
\qquad
q(1)=p(1)-\lambda=0,
\qquad
q(2)=p(2)-\lambda=0.
\]
Also hat $q$ die Nullstellen $0,1,2$.

\textbf{Schritt 2: Nullstellen als Faktoren schreiben.}

Wenn $q$ die Nullstellen $0,1,2$ hat, dann ist $q$ durch
\[
x(x-1)(x-2)
\]
teilbar.

Da $p\in \mathbb{R}_4[x]$ gilt, hat $p$ Grad h\"ochstens $4$.
Auch $q=p-\lambda$ hat Grad h\"ochstens $4$.
Der Faktor $x(x-1)(x-2)$ hat Grad $3$.
Also darf nach dem Herausziehen dieses Faktors nur noch ein Polynom vom Grad h\"ochstens $1$ \"ubrig bleiben:
\[
q(x)=(\alpha x+\beta)\,x(x-1)(x-2).
\]

Daher hat jedes $p\in W$ die Form
\[
p(x)=\lambda+\beta\,x(x-1)(x-2)+\alpha\,x^2(x-1)(x-2).
\]

\textbf{Schritt 3: Basis ablesen.}

Die drei freien Zahlen sind $\lambda,\beta,\alpha$.
Zu jeder freien Zahl geh\"ort ein Basis-Polynom:
\[
1,
\qquad
x(x-1)(x-2),
\qquad
x^2(x-1)(x-2).
\]

\[
\fbox{\parbox{0.93\textwidth}{\centering
Eine Basis von $W$ ist
\[
\{1,\; x(x-1)(x-2),\; x^2(x-1)(x-2)\}.
\]
Also ist $\dim(W)=3$.
}}
\]

\textbf{Kurze Kontrolle.}
Bei $x=0,1,2$ verschwinden die beiden Faktoren
\[
x(x-1)(x-2)
\quad\text{und}\quad
x^2(x-1)(x-2).
\]
Deshalb haben alle Polynome der Form
\[
\lambda+\beta x(x-1)(x-2)+\alpha x^2(x-1)(x-2)
\]
bei $0,1,2$ immer denselben Wert, n\"amlich $\lambda$.

\textbf{Warum sind die drei Polynome linear unabh\"angig?}
Ausmultipliziert sind
\[
x(x-1)(x-2)=x^3-3x^2+2x,
\]
\[
x^2(x-1)(x-2)=x^4-3x^3+2x^2.
\]
In einer Kombination
\[
A\cdot 1+B\cdot x(x-1)(x-2)+C\cdot x^2(x-1)(x-2)=0
\]
kommt der Term $x^4$ nur vom dritten Polynom. Also muss $C=0$ sein.
Dann kommt der Term $x^3$ nur noch vom zweiten Polynom. Also muss $B=0$ sein.
Dann bleibt $A\cdot 1=0$, also $A=0$.
Damit ist nur die triviale Kombination m\"oglich.

\bigskip
\hrule
\bigskip

%% ============================================================
\subsection*{$\square$ Beispiel 23(a)}
%% ============================================================

\textbf{Angabe.}
Gegeben ist das lineare Gleichungssystem in den Variablen $x,y,z,t,u$:
\[
\begin{cases}
x+2y+z+2t-2u=0,\\
3x+6y+z+2u=0,\\
2x+4y+z+t=0.
\end{cases}
\]
Gesucht sind eine Basis und die Dimension des L\"osungsraums.

\begin{begriffe}
\begin{itemize}[leftmargin=*, nosep]
    \item Der \textbf{L\"osungsraum} besteht aus allen Vektoren $(x,y,z,t,u)\in\mathbb{R}^5$, die alle drei Gleichungen erf\"ullen.
    \item Weil rechts immer $0$ steht, ist das System \textbf{homogen}. Deshalb ist die L\"osungsmenge automatisch ein Unterraum.
    \item Freie Variablen liefern die Parameter. Die Parametervektoren bilden die Basis des L\"osungsraums.
\end{itemize}
\end{begriffe}

\textbf{Schritt 1: Erweiterte Matrix aufstellen.}

\[
\left(\begin{array}{ccccc|c}
1 & 2 & 1 & 2 & -2 & 0 \\
3 & 6 & 1 & 0 & 2 & 0 \\
2 & 4 & 1 & 1 & 0 & 0
\end{array}\right)
\]

\textbf{Schritt 2: Gau\ss-Elimination.}

Wir benutzen die erste Zeile als Pivot-Zeile. Der erste Eintrag von $Z_1$ ist schon $1$,
also ist das ein guter erster Pivot.
Jetzt machen wir darunter in der $x$-Spalte Nullen.

\[
\begin{array}{rcl}
Z_2-3Z_1
&=&
(3,6,1,0,2,0)-3(1,2,1,2,-2,0) \\
&=&
(0,0,-2,-6,8,0),
\\[4pt]
Z_3-2Z_1
&=&
(2,4,1,1,0,0)-2(1,2,1,2,-2,0) \\
&=&
(0,0,-1,-3,4,0).
\end{array}
\]

Damit bekommen wir:
\[
\xrightarrow[\substack{Z_3\to Z_3-2Z_1}]{Z_2\to Z_2-3Z_1}
\left(\begin{array}{ccccc|c}
1 & 2 & 1 & 2 & -2 & 0 \\
0 & 0 & -2 & -6 & 8 & 0 \\
0 & 0 & -1 & -3 & 4 & 0
\end{array}\right).
\]

Die zweite Zeile ist genau das Doppelte der dritten Zeile.
Das hei\ss t: Diese beiden Zeilen sagen inhaltlich dasselbe.
Also bringt eine davon keine neue Information.
Wir machen deshalb $Z_2$ zur Nullzeile:
\[
Z_2-2Z_3
=
(0,0,-2,-6,8,0)-2(0,0,-1,-3,4,0)
=
(0,0,0,0,0,0).
\]

\[
\xrightarrow{Z_2\to Z_2-2Z_3}
\left(\begin{array}{ccccc|c}
1 & 2 & 1 & 2 & -2 & 0 \\
0 & 0 & 0 & 0 & 0 & 0 \\
0 & 0 & -1 & -3 & 4 & 0
\end{array}\right).
\]

Jetzt schieben wir die Nullzeile nach unten, weil Nullzeilen in der Stufenform immer unten stehen sollen.
Danach machen wir aus dem Pivot $-1$ eine $1$, indem wir die ganze Zeile mit $-1$ multiplizieren:
\[
\xrightarrow{Z_2\leftrightarrow Z_3}
\left(\begin{array}{ccccc|c}
1 & 2 & 1 & 2 & -2 & 0 \\
0 & 0 & -1 & -3 & 4 & 0 \\
0 & 0 & 0 & 0 & 0 & 0
\end{array}\right)
\xrightarrow{Z_2\to -Z_2}
\left(\begin{array}{ccccc|c}
1 & 2 & 1 & 2 & -2 & 0 \\
0 & 0 & 1 & 3 & -4 & 0 \\
0 & 0 & 0 & 0 & 0 & 0
\end{array}\right).
\]

Jetzt ist in der zweiten Zeile der Pivot in der $z$-Spalte.
\"Uber diesem Pivot steht in $Z_1$ noch eine $1$.
F\"ur die reduzierte Form soll dort eine $0$ stehen, also rechnen wir
\[
Z_1-Z_2
=
(1,2,1,2,-2,0)-(0,0,1,3,-4,0)
=
(1,2,0,-1,2,0).
\]

\[
\xrightarrow{Z_1\to Z_1-Z_2}
\left(\begin{array}{ccccc|c}
\piv{1} & 2 & 0 & -1 & 2 & 0 \\
0 & 0 & \piv{1} & 3 & -4 & 0 \\
0 & 0 & 0 & 0 & 0 & 0
\end{array}\right).
\]

\textbf{Fertige Version der Gau\ss-Elimination.}

Wenn man nur die fertige Matrix anschauen will, ist das die wichtige Endform:
\[
\boxed{
\left(\begin{array}{ccccc|c}
\piv{1} & 2 & 0 & -1 & 2 & 0 \\
0 & 0 & \piv{1} & 3 & -4 & 0 \\
0 & 0 & 0 & 0 & 0 & 0
\end{array}\right)
}
\]
Das ist die Form, bei der man mit der Gau\ss-Elimination fertig ist:
Die Pivot-Spalten sind sichtbar, und aus den zwei Nichtnull-Zeilen k\"onnen wir direkt die Gleichungen f\"ur $x$ und $z$ ablesen.

\textbf{Schritt 3: Pivot- und freie Variablen ablesen.}

Jetzt schauen wir nicht mehr auf die Rechenschritte, sondern nur auf die fertige Matrix.
Wichtig ist: Die Spalten geh\"oren der Reihe nach zu den Variablen
\[
x,\quad y,\quad z,\quad t,\quad u.
\]

Mit Spaltennamen sieht die fertige Matrix so aus:
\[
\begin{array}{c|ccccc|c}
 & x & y & z & t & u & 0 \\ \hline
Z_1 & \piv{1} & 2 & 0 & -1 & 2 & 0 \\
Z_2 & 0 & 0 & \piv{1} & 3 & -4 & 0 \\
Z_3 & 0 & 0 & 0 & 0 & 0 & 0
\end{array}
\]

Die Pivot-Spalten sind genau die Spalten, in denen ein markierter Pivot steht.
Hier stehen Pivots in der $x$-Spalte und in der $z$-Spalte.
Deshalb werden $x$ und $z$ sp\"ater ausgerechnet.

In den Spalten $y$, $t$ und $u$ steht kein Pivot.
Diese Variablen d\"urfen wir frei w\"ahlen.
Die freien Variablen sind also
\[
y,\qquad t,\qquad u.
\]

Jetzt ziehen wir die Gleichungen aus den Nichtnull-Zeilen der Matrix heraus.
Jede Zeile bedeutet:
\[
\text{Eintrag in Spalte }x\cdot x
+\text{Eintrag in Spalte }y\cdot y
+\cdots
=\text{rechte Seite}.
\]

\textbf{Zeile 1:}
\[
1x+2y+0z-1t+2u=0.
\]
Den Term $0z$ kann man weglassen:
\[
x+2y-t+2u=0.
\]

\textbf{Zeile 2:}
\[
0x+0y+1z+3t-4u=0.
\]
Die Nullterme kann man weglassen:
\[
z+3t-4u=0.
\]

Also liefert die Matrix insgesamt die zwei Gleichungen
\[
x+2y-t+2u=0,
\qquad
z+3t-4u=0.
\]

Jetzt l\"osen wir nach den Pivot-Variablen $x$ und $z$ auf.
Alles ohne Pivot bleibt rechts als freie Variable stehen:
\[
x=-2y+t-2u,
\qquad
z=-3t+4u.
\]

\textbf{Schritt 4: Freie Variablen als Parameter setzen.}

Setze
\[
y=s,\qquad t=r,\qquad u=q
\quad (s,r,q\in\mathbb{R}).
\]
Dann ist
\[
x=-2s+r-2q,
\qquad
z=-3r+4q.
\]

Damit sieht jeder L\"osungsvektor so aus:
\begin{align*}
(x,y,z,t,u)
&=(-2s+r-2q,\; s,\; -3r+4q,\; r,\; q) \\
&= s(-2,1,0,0,0)
 + r(1,0,-3,1,0)
 + q(-2,0,4,0,1).
\end{align*}

\[
\fbox{\parbox{0.93\textwidth}{\centering
Eine Basis des L\"osungsraums ist
\[
\{(-2,1,0,0,0),\; (1,0,-3,1,0),\; (-2,0,4,0,1)\}.
\]
Die Dimension des L\"osungsraums ist $3$.
}}
\]

\bigskip
\hrule
\bigskip

%% ============================================================
\subsection*{$\square$ Beispiel 23(b)}
%% ============================================================

\textbf{Angabe.}
Mit dem Ergebnis aus 23(a) und dem passenden Satz sollen die Dimensionen von
\[
U=\operatorname{span}\big((1,2,1,2,-2),\; (3,6,1,0,2),\; (2,4,1,1,0)\big)
\]
und
\[
W=\operatorname{span}\big((1,3,2),\; (2,6,4),\; (1,1,1),\; (2,0,1),\; (-2,2,0)\big)
\]
bestimmt werden, ohne weitere Rechnungen.

\begin{wastun}
\begin{enumerate}[leftmargin=*, nosep]
    \item Die Aufgabe sagt: Wir sollen \textbf{nicht nochmal Gau\ss rechnen}.
    \item Stattdessen sollen wir das Ergebnis aus 23(a) verwenden.
    \item Daf\"ur bauen wir eine Matrix $A$, deren Zeilen und Spalten genau zu den Vektoren in $U$ und $W$ passen.
    \item Dann benutzen wir den Dimensionssatz: Kern-Dimension plus Rang ergibt Anzahl der Spalten.
\end{enumerate}
\end{wastun}

\textbf{Let}
\[
A=
\begin{pmatrix}
1 & 2 & 1 & 2 & -2 \\
3 & 6 & 1 & 0 & 2 \\
2 & 4 & 1 & 1 & 0
\end{pmatrix}.
\]

\begin{begriffe}
\begin{itemize}[leftmargin=*, nosep]
    \item \textbf{span} liest man ``Spann'' oder ``Span''. Das ist kein neues Rechenzeichen wie $+$ oder $\cdot$, sondern hei\ss t: \textbf{alle Linearkombinationen} der angegebenen Vektoren.
    \item Wenn dort $\operatorname{span}(v_1,v_2,v_3)$ steht, meint das:
    \[
    \{\lambda_1v_1+\lambda_2v_2+\lambda_3v_3 : \lambda_1,\lambda_2,\lambda_3\in\mathbb{R}\}.
    \]
    Also: Wir d\"urfen die Vektoren beliebig mit Zahlen multiplizieren und addieren.
    \item $U$ ist ein Unterraum von $\mathbb{R}^5$, weil die drei Vektoren in $U$ jeweils $5$ Eintr\"age haben:
    \[
    (1,2,1,2,-2),\quad (3,6,1,0,2),\quad (2,4,1,1,0).
    \]
    \item $W$ ist ein Unterraum von $\mathbb{R}^3$, weil die f\"unf Vektoren in $W$ jeweils $3$ Eintr\"age haben:
    \[
    (1,3,2),\quad (2,6,4),\quad (1,1,1),\quad (2,0,1),\quad (-2,2,0).
    \]
    \item Der Unterschied zwischen $U$ und $W$ ist also: $U$ wird von \textbf{3 langen Zeilenvektoren} in $\mathbb{R}^5$ aufgespannt, $W$ von \textbf{5 kurzen Spaltenvektoren} in $\mathbb{R}^3$.
    \item Die Zeilen von $A$ sind genau die drei Vektoren, die $U$ aufspannen. Deshalb ist $U$ der \textbf{Zeilenraum} von $A$.
    \item Die Spalten von $A$ sind genau die f\"unf Vektoren, die $W$ aufspannen. Deshalb ist $W$ der \textbf{Spaltenraum} von $A$.
    \item Der L\"osungsraum aus 23(a) ist der Kern von $A$, also $\ker(A)$, denn in 23(a) haben wir genau $A\cdot(x,y,z,t,u)^T=0$ gel\"ost.
\end{itemize}
\end{begriffe}

\textbf{Satz.}
F\"ur eine Matrix $A\in\mathbb{R}^{m\times n}$ gilt der Dimensionssatz
\[
\dim(\ker A)+\operatorname{rang}(A)=n.
\]
Dabei ist $n$ die Anzahl der Spalten von $A$.
Au\ss erdem gilt:
\[
\operatorname{rang}(A)
= \dim(\text{Spaltenraum von }A)
= \dim(\text{Zeilenraum von }A).
\]

\begin{begriffe}[title={Was bedeuten die Begriffe im Satz?}]
\begin{itemize}[leftmargin=*, nosep]
    \item $\dim$ ist die Abk\"urzung f\"ur \textbf{Dimension}. Die Dimension sagt: Wie viele Vektoren stehen in einer Basis? Beispiel: Wenn ein Raum eine Basis mit $3$ Vektoren hat, dann hat er Dimension $3$.
    \item $\ker(A)$ hei\ss t \textbf{Kern} von $A$. Das ist die L\"osungsmenge von
    \[
    A\vec{x}=0.
    \]
    Genau diesen Kern haben wir in 23(a) berechnet.
    \item $\dim(\ker A)$ hei\ss t also: Wie viele freie Parameter hat die L\"osungsmenge von $A\vec{x}=0$? In 23(a) waren das $3$ freie Parameter, also $\dim(\ker A)=3$.
    \item $\operatorname{rang}(A)$ hei\ss t \textbf{Rang} von $A$. Praktisch bedeutet das: Wie viele Pivot-Spalten bzw. wie viele echte unabh\"angige Informationen hat die Matrix?
    \item Der \textbf{Spaltenraum} von $A$ ist der Span der Spalten von $A$. Seine Dimension ist die Anzahl der unabh\"angigen Spalten.
    \item Der \textbf{Zeilenraum} von $A$ ist der Span der Zeilen von $A$. Seine Dimension ist die Anzahl der unabh\"angigen Zeilen.
    \item Der wichtige Satz sagt: Zeilenraum und Spaltenraum haben immer dieselbe Dimension. Diese gemeinsame Dimension nennt man den Rang.
\end{itemize}
\end{begriffe}

\textbf{Schritt 1: Ergebnis aus 23(a) einsetzen.}

In 23(a) haben wir gefunden:
\[
\dim(\ker A)=3.
\]
Die Matrix $A$ hat $5$ Spalten. Also ist $n=5$.

Mit dem Dimensionssatz:
\[
\dim(\ker A)+\operatorname{rang}(A)=5.
\]
Einsetzen:
\[
3+\operatorname{rang}(A)=5.
\]
Also:
\[
\operatorname{rang}(A)=2.
\]

\textbf{Schritt 2: Dimension von $U$.}

$U$ ist der Zeilenraum von $A$, weil $U$ von den Zeilen von $A$ aufgespannt wird.
Die Dimension des Zeilenraums ist der Rang:
\[
\dim(U)=\operatorname{rang}(A)=2.
\]

\textbf{Schritt 3: Dimension von $W$.}

$W$ ist der Spaltenraum von $A$, weil $W$ von den Spalten von $A$ aufgespannt wird.
Die Dimension des Spaltenraums ist ebenfalls der Rang:
\[
\dim(W)=\operatorname{rang}(A)=2.
\]

\[
\fbox{\parbox{0.93\textwidth}{\centering
\[
\dim(U)=2,
\qquad
\dim(W)=2.
\]
}}
\]

\bigskip
\hrule
\bigskip

%% ============================================================
\subsection*{$\square$ Beispiel 24}
%% ============================================================

\textbf{Angabe.}
\[
A=
\begin{pmatrix}
6 & 1 & 3 \\
2 & 0 & 1 \\
1 & 2 & 0
\end{pmatrix}
\in \mathbb{R}^{3\times 3}.
\]

\begin{itemize}[leftmargin=*]
    \item[(a)] Berechne $A^{-1}$ mit dem Algorithmus $(A\mid I)\to(I\mid A^{-1})$.
    \item[(b)] Verwende das Ergebnis aus (a), um
    \[
    A
    \begin{pmatrix}
    x\\y\\z
    \end{pmatrix}
    =
    \begin{pmatrix}
    11\\4\\-1
    \end{pmatrix}
    \]
    zu l\"osen.
\end{itemize}

\begin{begriffe}
\begin{itemize}[leftmargin=*, nosep]
    \item $A^{-1}$ ist die inverse Matrix. Sie macht die Wirkung von $A$ wieder r\"uckg\"angig.
    \item Der Algorithmus $(A\mid I)\to(I\mid A^{-1})$ bedeutet: Wir schreiben rechts neben $A$ die Einheitsmatrix und machen links per Zeilenumformungen die Einheitsmatrix. Was rechts entsteht, ist $A^{-1}$.
    \item Wenn $A\vec{x}=\vec{b}$ gilt, dann ist $\vec{x}=A^{-1}\vec{b}$.
\end{itemize}
\end{begriffe}

%% --- 24(a) ---
\subsubsection*{24(a) Inverse berechnen}

\textbf{Schritt 1: Erweiterte Matrix aufstellen.}

\[
\left(\begin{array}{ccc|ccc}
6 & 1 & 3 & 1 & 0 & 0 \\
2 & 0 & 1 & 0 & 1 & 0 \\
1 & 2 & 0 & 0 & 0 & 1
\end{array}\right)
\]

\textbf{Schritt 2: Eine $1$ als ersten Pivot nach oben holen.}

Die dritte Zeile beginnt mit $1$. Das ist ein angenehmer Pivot.
Wir tauschen daher $Z_1$ und $Z_3$:
\[
\xrightarrow{Z_1\leftrightarrow Z_3}
\left(\begin{array}{ccc|ccc}
1 & 2 & 0 & 0 & 0 & 1 \\
2 & 0 & 1 & 0 & 1 & 0 \\
6 & 1 & 3 & 1 & 0 & 0
\end{array}\right).
\]

\textbf{Schritt 3: Unter dem ersten Pivot Nullen machen.}

\[
\xrightarrow[\substack{Z_3\to Z_3-6Z_1}]{Z_2\to Z_2-2Z_1}
\left(\begin{array}{ccc|ccc}
1 & 2 & 0 & 0 & 0 & 1 \\
0 & -4 & 1 & 0 & 1 & -2 \\
0 & -11 & 3 & 1 & 0 & -6
\end{array}\right).
\]

\textbf{Schritt 4: Zweiten Pivot zu $1$ machen.}

\[
\xrightarrow{Z_2\to -\frac14 Z_2}
\left(\begin{array}{ccc|ccc}
1 & 2 & 0 & 0 & 0 & 1 \\
0 & 1 & -\frac14 & 0 & -\frac14 & \frac12 \\
0 & -11 & 3 & 1 & 0 & -6
\end{array}\right).
\]

\textbf{Schritt 5: In der zweiten Spalte oben und unten Nullen machen.}

\[
\xrightarrow[\substack{Z_3\to Z_3+11Z_2}]{Z_1\to Z_1-2Z_2}
\left(\begin{array}{ccc|ccc}
1 & 0 & \frac12 & 0 & \frac12 & 0 \\
0 & 1 & -\frac14 & 0 & -\frac14 & \frac12 \\
0 & 0 & \frac14 & 1 & -\frac{11}{4} & -\frac12
\end{array}\right).
\]

\textbf{Schritt 6: Dritten Pivot zu $1$ machen.}

\[
\xrightarrow{Z_3\to 4Z_3}
\left(\begin{array}{ccc|ccc}
1 & 0 & \frac12 & 0 & \frac12 & 0 \\
0 & 1 & -\frac14 & 0 & -\frac14 & \frac12 \\
0 & 0 & 1 & 4 & -11 & -2
\end{array}\right).
\]

\textbf{Schritt 7: In der dritten Spalte oben Nullen machen.}

\[
\xrightarrow[\substack{Z_2\to Z_2+\frac14 Z_3}]{Z_1\to Z_1-\frac12 Z_3}
\left(\begin{array}{ccc|ccc}
1 & 0 & 0 & -2 & 6 & 1 \\
0 & 1 & 0 & 1 & -3 & 0 \\
0 & 0 & 1 & 4 & -11 & -2
\end{array}\right).
\]

Links steht jetzt $I$. Rechts steht also $A^{-1}$:
\[
\fbox{\parbox{0.93\textwidth}{\centering
\[
A^{-1}=
\begin{pmatrix}
-2 & 6 & 1 \\
1 & -3 & 0 \\
4 & -11 & -2
\end{pmatrix}.
\]
}}
\]

%% --- 24(b) ---
\subsubsection*{24(b) Gleichungssystem mit $A^{-1}$ l\"osen}

Wir wollen
\[
A
\begin{pmatrix}
x\\y\\z
\end{pmatrix}
=
\begin{pmatrix}
11\\4\\-1
\end{pmatrix}
\]
l\"osen.

Weil $A^{-1}$ die inverse Matrix ist, multiplizieren wir beide Seiten von links mit $A^{-1}$:
\[
A^{-1}A
\begin{pmatrix}
x\\y\\z
\end{pmatrix}
=
A^{-1}
\begin{pmatrix}
11\\4\\-1
\end{pmatrix}.
\]

Da $A^{-1}A=I$, bleibt links nur der gesuchte Vektor:
\[
\begin{pmatrix}
x\\y\\z
\end{pmatrix}
=
A^{-1}
\begin{pmatrix}
11\\4\\-1
\end{pmatrix}.
\]

Jetzt setzen wir $A^{-1}$ ein:
\[
\begin{pmatrix}
x\\y\\z
\end{pmatrix}
=
\begin{pmatrix}
-2 & 6 & 1 \\
1 & -3 & 0 \\
4 & -11 & -2
\end{pmatrix}
\begin{pmatrix}
11\\4\\-1
\end{pmatrix}.
\]

Eintrag f\"ur Eintrag:
\begin{align*}
x &= -2\cdot 11+6\cdot 4+1\cdot(-1)
   = -22+24-1
   = 1, \\
y &= 1\cdot 11+(-3)\cdot 4+0\cdot(-1)
   = 11-12+0
   = -1, \\
z &= 4\cdot 11+(-11)\cdot 4+(-2)\cdot(-1)
   = 44-44+2
   = 2.
\end{align*}

\[
\fbox{\parbox{0.93\textwidth}{\centering
\[
\begin{pmatrix}
x\\y\\z
\end{pmatrix}
=
\begin{pmatrix}
1\\-1\\2
\end{pmatrix}.
\]
Also ist $x=1$, $y=-1$, $z=2$.
}}
\]

\textbf{Probe.}
\[
A
\begin{pmatrix}
1\\-1\\2
\end{pmatrix}
=
\begin{pmatrix}
6\cdot1+1\cdot(-1)+3\cdot2 \\
2\cdot1+0\cdot(-1)+1\cdot2 \\
1\cdot1+2\cdot(-1)+0\cdot2
\end{pmatrix}
=
\begin{pmatrix}
11\\4\\-1
\end{pmatrix}.
\]
Die Probe passt.

\end{document}
